\documentclass[12pt]{article}
\usepackage{hyperref}

\usepackage{graphicx}

\usepackage{mathtools}
\DeclarePairedDelimiter\ceil{\lceil}{\rceil}
\DeclarePairedDelimiter\floor{\lfloor}{\rfloor}


\usepackage{amsmath}
\usepackage{amsfonts}
\usepackage{amssymb}

\usepackage{listings}
\lstset{language=C++,
                basicstyle=\ttfamily,
                keywordstyle=\color{blue}\ttfamily,
                stringstyle=\color{red}\ttfamily,
                commentstyle=\color{green}\ttfamily,
                morecomment=[l][\color{magenta}]{\#}
}

\usepackage{breqn}
\usepackage[]{algorithm2e}
\usepackage{physics}

\newcommand\fnurl[2]{%
  \href{#2}{#1}\footnote{\url{#2}}%
}

\usepackage{cancel}

\newcommand{\Four}{\mathcal{F}}
\renewcommand{\(}{\left(}
\renewcommand{\)}{\right)}
%% \renewcommand{\{}{\left{}
%% \renewcommand{\}}{\right}}


\title{Twiddle Table Notes for rocFFT}
\author{The rocFFT Team}



\def\no#1{\nobreak{#1}} % stop equation breaking in dmath

\begin{document}
\maketitle

\section{Introduction}

The twiddle table for a FFT of length $N$ are just the $N$ roots of unity:
\begin{dmath}
  \omega_N^k = e^{\frac{k 2 \pi i}{N}} = \cos\frac{k 2 \pi}{N}
  +i\sin\frac{k 2 \pi}{N} \text{, for } k\nobreak{=}0,\dots,N-1.
\end{dmath}
Figure~\ref{rooteg} shows these in the complex plane.
\begin{figure}[htbp]
  \centering
  \includegraphics{rootn}
  \caption{Roots of unity for $N=5$.}
  \label{rooteg}
\end{figure}


\subsection{Symmetries in Twiddle Tables}

\subsubsection{General $N$}
It's easy to see in Figure~\ref{rooteg} that there are a number of
symmtries that we can take advantage of.  For general, $N$, we know
that, for $k=N-1$,
\begin{dmath}
  e^{\frac{(N-1)2 \pi i}{N}}
  =  \cancelto{1}{e^{\frac{N2 \pi i}{N}}}\, e^{\frac{-2 \pi i}{N}}
  = \(e^{\frac{2 \pi i}{N}}\)^*.
\end{dmath}
That is, the $k=(N-1)$th root of unity is the complex conjugate of the
$k=1$ root of unity.  This can be generalised:
\begin{dmath}
  \omega_N^{N - k}
  = \omega_N^{-k}
  =  {\omega_N^{k}}^*
\end{dmath}
so, for $k=\floor{\frac{N}{2}} + 1,\dots,N-1$,
\begin{dmath}
  \omega_N^{k}
  = {\omega_N^{N-k}}^*
\end{dmath}
so we can compute the $k$th twiddle as the complex conjugate of the
$N-k$th twiddle.  This reduces the size of the twiddle table by a
factor of 2.

\subsubsection{Even $N$}
If $N$ is even, then we have another symmetry, as shown in
Figure~\ref{rooteven}: Keeping only half of the complex plane, we also
have that $\omega_{N}^{N/2-1} = -\(\omega_{N}^{1}\)^*$.
\begin{figure}[htbp]
  \centering
  \includegraphics{rooteven}
  \caption{Roots of unity for $N=14$.}
  \label{rooteven}
\end{figure}
In general, for even $N$, we have
\begin{dmath}
  \omega_N^{\frac{N}{2} - k}
  = \omega_N^{N/2}\omega_N^{-k}
  = -\(\omega_N^{k}\)^*
  = -\Re{\omega_N^{k}\}} + i \Im{\omega_N^{k}\}}.
 \end{dmath}
This allows us to reduce the twiddle table size by a further factor of
2, keeping only $1/4$ of the original values.  The computation is just
changing the sing of the real part.


\subsubsection{$N|4$}

If $N|4$, then there is another useful symmetry, which requires no
additional FLOPS:
\begin{figure}[htbp]
  \centering
  \includegraphics{rootdiv4}
  \caption{Roots of unity for $N=20$.}
  \label{rootdiv4}
\end{figure}
and we have
\begin{dmath}
  \omega_N^{\frac{N}{4} - k}
  = \omega_N^{\frac{N}{4}}\omega_N^{-k}
  = \omega_4\(\omega_N^{k}\)^*
  = i\(\omega_N^{k}\)^*
  = \Im{\omega_N^{k}} + i\Re{\omega_N^{k}}
\end{dmath}
In other words, we can reduce the twiddle table by a factor of 8 if
$N|4$.

\subsubsection{$N|8$}

If $N|8$, then there is a final symmetry to take advantage of which
involves the online computation of $\sqrt{2}$.
\begin{figure}[htbp]
  \centering
  \includegraphics{rootdiv8}
  \caption{Roots of unity for $N=32$.}
  \label{rootdiv8}
\end{figure}
Here, we have
\begin{dmath}
  \omega_N^{\frac{N}{8} - k}
  = \omega_N^{\frac{N}{8}}\omega_N^{-k}
  = \omega_8\(\omega_N^{k}\)^*
  = \(\sqrt{2} + i \sqrt{2}\) \(\omega_N^{k}\)^*.
\end{dmath}
This allows us to store only $1/16$ of the original twiddle table.

\end{document}
